%%%%%%%%%%%%%%%%%%%%%%%%%%%%%%%%%%%%%%%%%%%%%%%%%%
\needspace{12\baselineskip}
\section{Algorithmen und Programmiersprachen}
%%%%%%%%%%%%%%%%%%%%%%%%%%%%%%%%%%%%%%%%%%%%%%%%%%
\worum
Ein \textbf{Algorithmus} ist eine detaillierte und eindeutige Vorschrift zur schrittweisen Lösung eines Problems. Er ist die \emph{konzeptionelle} Lösung. Ein \textbf{Programm} ist die auf einem Rechner ausführbare Lösung, also der Algorithmus in einer Programmiersprache. Vor dem Algorithmus steht die \textbf{Spezifikation}: Sie legt fest, welche Eingaben erlaubt sind (Vorbedingung) und was für die Ergebnisse gelten muss (Nachbedingung).

Algorithmen lassen sich auf verschiedene Arten aufschreiben: in mathematischer Notation (oft rekursiv), umgangssprachlich, als \textbf{Flussdiagramm}, als Pseudocode oder direkt in einer Programmiersprache.

Damit ein Rechner ein Programm ausführen kann, muss es in Maschinenbefehle umgesetzt werden. Dafür gibt es zwei Wege: Ein \textbf{Compiler} übersetzt das ganze Programm \emph{vor} der Ausführung in Maschinencode. Ein \textbf{Interpreter} analysiert das Programm \emph{zur Laufzeit} Anweisung für Anweisung und führt es direkt aus. Java ist eine Mischform: Der Compiler (\code{javac}) erzeugt plattformunabhängigen Byte-Code, den die virtuelle Maschine (JVM) interpretiert.

\begin{merke}[Eigenschaften von Algorithmen]
\begin{itemize}[itemsep=0.1em]
\item \textbf{terminierend:} bricht für jede zulässige Eingabe nach endlich vielen Schritten ab.
\item \textbf{determiniert:} liefert bei gleichen Eingaben stets das gleiche Ergebnis.
\item \textbf{deterministisch:} zu jedem Zeitpunkt ist der nächste Schritt eindeutig bestimmt.
\item \textbf{partiell korrekt:} jedes berechnete Ergebnis entspricht der Spezifikation. \\
\textbf{total korrekt:} partiell korrekt und terminierend.
\item \textbf{effektiv:} erreicht das Ziel. \ \textbf{effizient:} erreicht es mit möglichst geringem Aufwand.
\end{itemize}
\end{merke}

\begin{merke}[Compiler und Interpreter]
\begin{tabular}{@{}L{7.2cm}L{7.2cm}@{}}
\textbf{Compiler} & \textbf{Interpreter} \\
übersetzt vor der Ausführung & übersetzt während der Ausführung \\
Schleifen werden nur einmal übersetzt, das Programm läuft schnell & Anweisungen werden evtl. mehrfach übersetzt, das Programm läuft langsamer \\
viele Fehler werden schon beim Übersetzen gefunden & Fehler zeigen sich erst zur Laufzeit \\
Beispiele: C, C++, Pascal & Beispiele: Python, PHP, Basic
\end{tabular}
\end{merke}

\beispiele
\beispiel{Der euklidische Algorithmus in drei Darstellungen}
Gesucht ist der größte gemeinsame Teiler zweier Zahlen $a \geq 0$ und $b \geq 0$.

\smallskip
\noindent\textbf{Mathematisch (rekursiv):}
\[
\mathit{ggT}(a, b) =
\begin{cases}
a & \text{falls } b = 0 \\
\mathit{ggT}(b,\ a \bmod b) & \text{sonst}
\end{cases}
\]
\noindent\begin{minipage}[t]{0.56\linewidth}
\textbf{Als Flussdiagramm:}

\medskip
\begin{tikzpicture}
	\node[startstop] (start) at (0,0) {Start};
	\node[io] (in) at (0,-1.2) {Eingabe $a$, $b$};
	\node[decision] (q) at (0,-2.7) {$b = 0$?};
	\node[process] (p1) at (0,-4.3) {$r = a \bmod b$ \\ $a = b$ \\ $b = r$};
	\node[io] (out) at (-3.6,-2.7) {Ausgabe $a$};
	\node[startstop] (end) at (-3.6,-4.0) {Ende};
	\draw[arrow] (start) -- (in);
	\draw[arrow] (in) -- (q);
	\draw[arrow] (q) -- node[lbl] {nein} (p1);
	\draw[arrow] (p1.east) -- (2.3,-4.3) -- (2.3,-2.7) -- (q.east);
	\draw[arrow] (q) -- node[lbl] {ja} (out);
	\draw[arrow] (out) -- (end);
\end{tikzpicture}
\end{minipage}\hfill
\begin{minipage}[t]{0.4\linewidth}
\textbf{Ablauf für $a = 960$, $b = 230$:}

\medskip
\begin{tabular}{l|r|r|r}
\textbf{Durchlauf} & $\boldsymbol{r}$ & $\boldsymbol{a}$ & $\boldsymbol{b}$ \\ \hline
Start & & 960 & 230 \\
1. & 40 & 230 & 40 \\
2. & 30 & 40 & 30 \\
3. & 10 & 30 & 10 \\
4. & 0 & 10 & 0
\end{tabular}

\medskip
Ausgabe: 10

\medskip
Eine solche Tabelle heißt \emph{Schreibtischtest}. Mit ihr prüft man einen Algorithmus von Hand.
\end{minipage}

\beispiel{Determiniert und deterministisch}
\begin{itemize}[itemsep=0.1em]
\item \enquote{Wähle eine beliebige Zahl zwischen 1 und 5, addiere sie zur Eingabe und gib \textbf{42} aus.} Der Ablauf ist nicht eindeutig festgelegt (nicht deterministisch), das Ergebnis ist aber immer dasselbe (determiniert).
\item \enquote{Wirf eine Münze. Bei Kopf gib die Eingabe aus, bei Zahl das Doppelte.} Dieser Algorithmus ist weder deterministisch noch determiniert.
\end{itemize}

\beispiel{Terminiert das?}
Die Collatz-Folge (Ulam-Algorithmus) startet mit einer beliebigen Zahl $n$. Ist $n$ gerade, wird halbiert, sonst wird $3n + 1$ gebildet, bis 1 erreicht ist: $6, 3, 10, 5, 16, 8, 4, 2, 1$. Bis heute ist nicht bewiesen, dass dieses Verfahren für \emph{jede} Startzahl endet.

\beispiel{Fehlerarten}
\begin{itemize}[itemsep=0.1em]
\item \textbf{Syntaxfehler:} Verstoß gegen die Grammatik der Sprache, z.\,B. ein vergessenes Semikolon. Der Compiler meldet ihn.
\item \textbf{Semantischer Fehler:} syntaktisch richtig, aber ohne gültige Bedeutung, z.\,B. \code{int i = 8.8;} (zur Übersetzungszeit erkennbar) oder ein Zugriff außerhalb eines Arrays (Laufzeitfehler).
\item \textbf{Denkfehler:} Das Programm läuft, liefert aber nicht das gewünschte Ergebnis.
\end{itemize}

\aufgaben
\aufgabe[\K]{Begriffe}
\begin{enumerate}[label=\alph*)]
\item Was versteht man unter einem Compiler? Worin unterscheidet er sich von einem Interpreter?
\item Wann wird ein Algorithmus als determiniert bezeichnet? Geben Sie ein Beispiel für einen nicht determinierten Algorithmus.
\item Erklären Sie den Unterschied zwischen partieller und totaler Korrektheit.
\item Warum bezeichnet man Java als plattformunabhängig? Welche Rolle spielen dabei Byte-Code und JVM?
\end{enumerate}

\aufgabe{Schreibtischtest}
Führen Sie den folgenden Algorithmus für die Eingabe $n = 4711$ aus. Notieren Sie die Werte von $n$ und $s$ zu Beginn und nach jedem Schleifendurchlauf in einer Tabelle. Was gibt der Algorithmus aus, und was berechnet er allgemein? (\emph{div} ist die ganzzahlige Division.)
\begin{center}
\begin{tikzpicture}
	\node[startstop] (start) at (0,0) {Start};
	\node[io] (in) at (3,0) {Eingabe $n$};
	\node[process, minimum width=1.6cm] (init) at (6,0) {$s = 0$};
	\node[decision] (q) at (9,0) {$n > 0$?};
	\node[process] (p1) at (9,-1.8) {$s = s + n \bmod 10$ \\ $n = n \;\mathrm{div}\; 10$};
	\node[io] (out) at (12.4,0) {Ausgabe $s$};
	\node[startstop] (end) at (12.4,-1.8) {Ende};
	\draw[arrow] (start) -- (in);
	\draw[arrow] (in) -- (init);
	\draw[arrow] (init) -- (q);
	\draw[arrow] (q) -- node[lbl] {ja} (p1);
	\draw[arrow] (p1.west) -- (6.6,-1.8) -- (6.6,-0.9) -- (7.5,-0.9) -- (q.south west);
	\draw[arrow] (q) -- node[lbl] {nein} (out);
	\draw[arrow] (out) -- (end);
\end{tikzpicture}
\end{center}

\aufgabe{Flussdiagramme entwerfen}
Entwerfen Sie jeweils einen Algorithmus als Flussdiagramm. Für die Ein- und Ausgabe dürfen Sie die Bausteine \emph{Eingabe} und \emph{Ausgabe} verwenden.
\begin{enumerate}[label=\alph*)]
\item Bestimmen Sie das Maximum dreier ganzer Zahlen $a$, $b$ und $c$.
\item Bestimmen Sie, ob $n \in \mathbb{N}_0$ gerade ist. Verwenden Sie nur den Minusoperator und Vergleiche.
\item Berechnen Sie den Quotienten und den Rest der ganzzahligen Division $x : y$ für $x \in \mathbb{N}_0$, $y \in \mathbb{N}$. Erlaubt sind nur Vergleichs-, Plus- und Minusoperator.
\item Berechnen Sie die Potenz $x^y$ für $x, y \in \mathbb{N}_0$. Erlaubt sind Vergleiche, Plus, Minus und die Multiplikation.
\end{enumerate}

\aufgabe{Rekursive Spezifikation}
Die Addition lässt sich allein mit dem Nachfolger ($+1$) und dem Vorgänger ($-1$) rekursiv definieren:
\[
\mathit{sum}(a, b) =
\begin{cases}
a & \text{falls } b = 0 \\
\mathit{sum}(a,\ b - 1) + 1 & \text{falls } b > 0
\end{cases}
\]
\begin{enumerate}[label=\alph*)]
\item Werten Sie $\mathit{sum}(4, 3)$ Schritt für Schritt aus.
\item Geben Sie in gleicher Weise eine rekursive Funktion $\mathit{multi}(x, y)$ an, die das Produkt zweier Zahlen aus $\mathbb{N}_0$ nur mit Addition und Subtraktion berechnet.
\item Geben Sie eine rekursive Funktion $\mathit{rest}(a, b)$ für $a \bmod b$ an ($a \geq 0$, $b > 0$).
\end{enumerate}

\aufgabe{Eigenschaften prüfen}
Der folgende Test soll feststellen, ob eine Zahl $z > 2$ eine Primzahl ist.
\begin{enumerate}[itemsep=0.1em, label=(\arabic*)]
\item Wähle ein beliebiges, noch nicht getestetes $n$ mit $1 < n < z$.
\item Falls $n$ ein Teiler von $z$ ist: Ausgabe \enquote{Keine Primzahl}, gehe zu (5).
\item Falls alle $n$ getestet sind: Ausgabe \enquote{Primzahl}, gehe zu (5).
\item Gehe zu (1).
\item Fertig.
\end{enumerate}
Ist der Algorithmus determiniert? Ist er deterministisch? Terminiert er? Ist er effizient? Begründen Sie jeweils kurz.

\aufgabe{Fehler einordnen}
Um welche Fehlerart handelt es sich jeweils (Syntaxfehler, semantischer Fehler zur Übersetzungszeit, Laufzeitfehler, Denkfehler)?
\begin{enumerate}[label=\alph*)]
\item \code{System.out.println("Hallo")} \quad (ohne Semikolon)
\item \code{double mittel = (a + b) / 2;} \quad mit \code{int a = 3, b = 4}, erwartet wird 3.5
\item \code{int[] f = new int[5]; f[5] = 1;}
\item \code{int zahl; int zahl2 = zahl;}
\end{enumerate}

\aufgabe[\S]{Primzahlen bis $n$}
Formulieren Sie einen Algorithmus als Flussdiagramm, der alle Primzahlen von 2 bis $n$ ausgibt. Erlaubt sind die Vergleichsoperatoren sowie $+$, $-$ und $\bmod$. Überlegen Sie anschließend: Bis zu welchem Teiler muss man höchstens prüfen?

\ueben
\aufgabe{Schreibtischtest ggT}
Führen Sie den euklidischen Algorithmus aus dem ersten Beispiel für die Eingaben $(84, 36)$, $(17, 5)$ und $(100, 75)$ aus. Notieren Sie die Werte von $r$, $a$ und $b$ nach jedem Durchlauf.

\aufgabe{Pseudocode lesen}
Was gibt der folgende Algorithmus für die Eingaben $n = 40$ und $n = 1000$ aus? Was berechnet er allgemein?
\begin{konsole}[Pseudocode]
lies n
z = 0
solange n > 1
    n = n div 2
    z = z + 1
ausgabe z
\end{konsole}

\aufgabe{Einfache Flussdiagramme}
Entwerfen Sie jeweils ein Flussdiagramm.
\begin{enumerate}[label=\alph*)]
\item Ausgabe des Betrags einer ganzen Zahl
\item Summe der Zahlen von 1 bis $n$
\item Fakultät $n!$ für $n \in \mathbb{N}_0$
\item Countdown von $n$ bis 1, danach die Ausgabe \enquote{Start}
\item Anzahl der Ziffern einer natürlichen Zahl
\end{enumerate}

\aufgabe{Rekursive Definitionen auswerten}
Werten Sie mit den Definitionen aus dem Kapitel Schritt für Schritt aus.
\begin{multicols}{3}
\begin{enumerate}[label=\alph*)]
\item $\mathit{ggT}(48, 18)$
\item $\mathit{sum}(2, 4)$
\item $\mathit{ggT}(7, 0)$
\end{enumerate}
\end{multicols}

\aufgabe{Eigenschaften zuordnen}
Welche der Eigenschaften \emph{terminierend}, \emph{determiniert} und \emph{deterministisch} treffen jeweils zu?
\begin{enumerate}[label=\alph*)]
\item \enquote{Lies $n$. Solange $n \neq 0$: verringere $n$ um 2. Gib \emph{fertig} aus.} (für $n \in \mathbb{N}_0$)
\item \enquote{Würfle. Gib die Augenzahl aus.}
\item \enquote{Lies eine Liste. Wähle zufällig zwei Elemente und vertausche sie, falls sie in der falschen Reihenfolge stehen. Wiederhole, bis die Liste sortiert ist. Gib die Liste aus.}
\item der euklidische Algorithmus
\end{enumerate}

\aufgabe{Compiler oder Interpreter?}
Ordnen Sie die Aussagen zu.
\begin{enumerate}[label=\alph*)]
\item Das Programm wird vor der Ausführung vollständig übersetzt.
\item Ein Tippfehler in Zeile 500 fällt erst auf, wenn diese Zeile erreicht wird.
\item Für jede Plattform wird ein eigenes ausführbares Programm erzeugt.
\item Eine Schleife mit 1000 Durchläufen wird nur einmal übersetzt.
\item Wie passt Java in dieses Schema?
\end{enumerate}
